% =====================================================================================================
\section{Discussion}\label{sec:discussion}

\paragraph{Strategy, persuasion and discrimination.}
One-step lookahead increased the primary judge's final guilt score under the advocate objective. The increase was
larger for guilty defendants, but it also affected innocent defendants and raised wrongful convictions under the
threshold verdict rule. Strategic selection can therefore improve discrimination while also increasing the cost
of false accusations. The estimated benefit was smaller under held-out judges \StudyTwoTransferText{}
This evaluator dependence is consistent with specialisation to a proxy measure \citep{manheim2018,gao2023},
although it does not by itself distinguish every source of disagreement between judges.

A lookahead defence reduced guilt scores and convictions. In the estimated zero-sum matrix based on final guilt
scores, both players preferred lookahead, and the resulting equilibrium score was close to that under random
selection on both sides. This result concerns the tested policy set and probability surrogate. In particular, the
defence's lookahead includes access to an internal judge that receives withheld evidence; its gain cannot be
attributed solely to reasoning over the defendant's own file.

\paragraph{Information and evaluation.}
The evidence-disclosure analysis associates early confrontation with better discrimination, and the branched
question experiment finds that outcome-informed selection improves discrimination more than the minister's
simulated scores do. These observations suggest a limitation of the planner's respondent model: hypothetical
suspects lack the private accounts that constrain the actual simulated defendants. Testing that explanation
would require varying the respondent model while holding the other policy components fixed. Changing the goal
instruction and score function alone did not produce a clear improvement in truth-tracking.

The judges also showed substantial sensitivity to continued examination and to their instructions. Both guilty
and innocent defendants received higher guilt scores as questioning continued, and a guilt-leaning instruction
more than tripled the conviction rate on the same transcripts. A policy effect is consequently inseparable from
the evaluation rule under which it is measured. The stopping-rule proposition permits comparisons of those rules
with fixed policies, but does not predict how speakers would adapt if the evaluator were changed and disclosed.

\paragraph{Implications for the Turing test.}
Relabelling the types in \cref{ex:turing} connects these objectives to interactive AI detection: an advocate-like
examiner seeks a machine verdict, while a minister-like examiner seeks correct classification. The simulations
suggest questions for such tests rather than direct estimates of their performance. Does an interrogation policy
raise detection at the expense of false accusations? How much of its advantage persists under another detector?
Can verifiable information distinguish the two types more reliably than stylistic cues? The answers depend on
the information available to each participant and on how the evaluator responds to adversarial dialogue.

\paragraph{Limitations.}
All speakers and judges are LLMs. The estimated effects concern the specified simulator, and transfer to human
suspects, judges or Turing-test participants remains an empirical question. The 48 main cases were generated and
audited under one stimulus procedure; the case-level intervals assume a comparable synthetic case population.
All players use one model, and the policies consider only four or eight candidates with shallow search. The
additional model calls required by introspection and lookahead are part of the treatment, so the experiments do
not estimate performance at equal inference cost. The defence selector's access to an evidence-informed judge
also limits claims about what a defendant could achieve using only its own information.

Judge scores are uncalibrated, differ in scale, and for two models use a truncated-token probability summary.
Only two continuations per candidate were collected in Study 2. Three decision points failed, and their exclusion
may affect generalisation beyond completed points. The crossover intervals condition on estimated normalisation
scales, and equilibrium statements concern an estimated surrogate-payoff matrix. Hosted-model variation,
provider fallback and incomplete sharing of cached responses add uncertainty that the recorded case-level
intervals need not capture fully. Finally, the analysis plan was recorded internally rather than deposited in an
external registry; its hashes establish file identity but do not independently verify the stated timing.


% =====================================================================================================
\section{Conclusion}
Verdict games describe strategic exchanges whose outcomes are determined by a non-strategic classifier. Their
information structure separates the evidence available to each speaker from that available to the judge, and a
fixed-policy replay argument specifies when stopped verdicts can be recovered from complete transcripts.
In the simulated examinations, lookahead improved the advocate's score, partly improved discrimination and also
increased perceived guilt for innocent defendants. Its advantage depended on the evaluator and was counteracted
by a strategic defence. Effects of individual question choices diminished under subsequent baseline play.
These findings link the value of strategic language to the speaker's objective, its information and the judge
it seeks to influence.
